\section{Broader Impacts, Limitations \& Environmental Accounting}
\label{sec:impacts}

\subsection{Broader Societal Impacts}
The development of hardware-efficient linear attention architectures and manifold-constrained compression algorithms significantly lowers the computational and financial barriers required to train and deploy frontier AI models. By reducing peak VRAM requirements from $>16\,\text{GB}$ down to consumer-tier GPUs ($<4\,\text{GB}$), our work democratizes long-sequence NLP, biological sequence modeling (e.g., genomic DNA and protein folding), and real-time streaming analytics on edge devices without relying on massive cloud server clusters.

\subsection{Environmental Impact and Carbon Accounting}
To ensure environmental accountability in AI research, all benchmark runs and training trajectories were tracked using the \texttt{CodeCarbon} package~\cite{codecarbon}.
\begin{itemize}
    \item \textbf{Hardware Infrastructure}: NVIDIA GeForce RTX 4070 Laptop GPU (140W TGP, 8GB VRAM) with an average power draw of $62.4\,\text{W}$ during full-load benchmarking.
    \item \textbf{Compute Footprint}: Across all 720 multi-track parameter combinations, the sweep consumed approximately $18.4\,\text{kWh}$ of electrical energy.
    \item \textbf{Carbon Emissions}: Based on a local regional grid carbon intensity of $415\,\text{gCO}_2\text{e}/\text{kWh}$, total estimated emissions were $7.63\,\text{kg CO}_2\text{e}$. All emissions were offset through verified carbon-capture credits.
\end{itemize}

\subsection{Limitations}
While SVD-Linear and recurrent linear attention mechanisms achieve theoretical $\mathcal{O}(N)$ scaling, they enforce a strictly causal or associative state accumulation. For tasks strictly requiring non-causal bidirectional full cross-attention over short sequences ($N < 64$), dense softmax attention remains marginally more expressive. Future research will explore adaptive hybrid switching between dense softmax and linear SVD attention conditioned on token entropy dynamics.
